\section{Shun Poison, Seek Nectar}

As often happens with saints, legends surround them. Ashtavakra's name derives from the eight
deformities in his body. The cause of the deformities is unclear. Said to be enlightened already in the womb,
Ashtavakra was upset about the scriptures being recited improperly; he either squirmed or was cursed by his father,
resulting in his physical defects. More likely, the eight deformities relate to the eightfold path of Yoga. However, it
is clear that most men can only relate to the physical meaning of things and miss subtle connections.

In the Ashtavakra Gita, the Sage \textbf{Ashtavakra} is teaching the King \textbf{Janaka}. The King is the Will of the
people, hence must himself be detached from personal interests. So he asks Ashtavakra:

\begin{quotex}
How can knowledge be acquired? How can liberation be attained? How can renunciation come about? (1) 

\end{quotex}
Obviously, like mostly everyone, the King is expecting to be told a secret. He wants some intellectual knowledge that he
can repeat. Instead, Ashtavakra's answer comes from an unexpected direction:

\begin{quotex}
My child, if you are seeking liberation, shun the objects of the senses like poison; and seek

\begin{itemize}
\item forgiveness 
\item sincerity 
\item kindness
\item happiness 
\item truth 
\end{itemize}
as you would seek nectar. (2) 

\end{quotex}
These are the flowers of the meadow, the ``stages" to follow to slough off attachments and realize
one's true identity. Ashtavakra continues:

\begin{quotex}
You are neither earth, nor water, nor fire, nor air, nor ether. You are the witness of those five elements as
Consciousness. Understanding this is liberation. (3) 

\end{quotex}
These five elements are the primordial causes of physical manifestation. Hence, Janaka must realize that his true
identity transcends the physical world. It is the Observer, the unmoved mover, the Atman. Ashtavakra describes this:

\begin{quotex}
The Atman is the sole witness, all pervading, perfect, free Consciousness: actionless, unattached, desireless, at peace
with itself. It is only through an illusion that it appears to be involved with samsara [the phenomenal world]. (12) 

\end{quotex}
The King must recognize what is unchanging and constantly present in every act of consciousness, every experience of the
phenomenal world. When identified with the contents of consciousness, Janaka is deluded. Ashtavakra explains:

\begin{quotex}
It is through ignorance of the Atman that the phenomenal universe appears to be real, and this illusion disappears with
the realization of one's true nature. (27) 

\end{quotex}
Nevertheless, Janaka prefers poison to sweet nectar. Ashtavakra questions his attachments in sutras such as these:

\begin{quotex}
It is indeed strange that one abiding in the supreme transcendent non-duality and intent on liberation should be subject
to lust and \emph{weakened by amorous activities}. It is a strange fact of this world that a man, physically weak and
obviously at the end of his life, should lust for sensual pleasure even after being aware that lust is an enemy of
knowledge. (51-52) 

\end{quotex}
It must be difficult for Janaka, because man prefers pleasure and avoids pain. Yet, both are forms of attachment and the
Sage transcends duality itself:

\begin{quotex}
It means bondage when the mind desires something or grieves at something, rejects or accepts anything, feels happy or
angry with anything. It means liberation when the mind does not desire, nor grieve, nor reject, nor accept, nor feel
happy or angry. (79-80) 

\end{quotex}
The Sage must not be attached to different points of view:

\begin{quotex}
Who will not attain tranquility who, seeing the diversity of opinions among the many seers, saints and yogis, becomes
totally indifferent? (87) 

\end{quotex}
Following some initiatic path is not some automatic process. Some men grasp enlightenment intuitively; others may spend
a lifetime seeking, but not finding. The latter most likely because the ``seeker" is playing at it, adopting outward
forms and ignoring the inner transformation.

\begin{quotex}
The person with a keen intellect becomes enlightened even when the instruction is imparted casually, whereas without it
the immature seeker continues to remain confused even after a lifetime of seeking. Absence of attachment to
sense-objects is liberation; passion for sense-objects is bondage. Understand this fact, and then do as you please.
(126-127) 

\end{quotex}
Of course, what that Sage is pleased to do is not what the ignorant finds pleasing. The ignorant man is attached to the
forces of \emph{eros} and \emph{thumos}, he enjoys the intense and complex dramas resulting therefrom; moreover, he
regards them as the essence of being ``human". Liberation from them appears to him as bondage, while bondage to the
senses is regarded as ``freedom" to do as one pleases. Hence, he finds the spiritual life unattractive.

\begin{quotex}
Apperception of this Truth seems to render an eloquent, wise, and active person mute, dull and inactive. Knowledge of
Truth does not therefore appear attractive to those who still want to enjoy the pleasures of this world. (128) 

\end{quotex}
Many will boast of ``love" or having a ``good heart" and so on; they may even make great efforts to appear as such to
others. It feels good for a moment, but that is not the way of the Sage.

\begin{quotex}
The non-volitional, spontaneous, unrestricted behavior of the wise man is transparently open and sincere but not the
affected tranquility exhibited by the one who is still governed by personal motives and considerations. (228) 

\end{quotex}


\flrightit{Posted on 2012-10-17 by Cologero }

\begin{center}* * *\end{center}

\begin{footnotesize}\begin{sffamily}



\texttt{Master Dogen on 2012-10-19 at 01:28 said: }

The faux zenned-out state — often characterized by a creepy soft voice and a slow yet somehow
aggressive smile — seems to me to be even farther from the true path than is a state of erotic
dissipation. 

But, what about when Ashtavakra's words seem correct, and yet still cannot cut through the veil?
That is, the ears are not deaf, but the heart is? I suppose each spirit will reveal itself when it wants to. But before
one wants to, one can WANT to want, though that's a rather anguished position to put oneself in.


\hfill

\texttt{Cologero on 2012-10-19 at 20:05 said: }

``the ears are not deaf, but the heart is"

Who can believe it? Who can believe that the end of desire is the end of suffering, when all our experience shows us
that the end of suffering is the satisfaction of desire?


\end{sffamily}\end{footnotesize}
